\subsection{Rancangan Model PBAC pada DecMed}
\label{subsec:rancangan-model-pbac-decmed}

Subbab ini merumuskan \textit{policy} akses atas objek. Berdasarkan analisis pada Subbab \ref{subsec:analisis-pendekatan-model-pbac}, elemen dan relasi \textit{Policy Machine} diadaptasi untuk memodelkan hubungan antara pengguna, atribut pengguna, segmen RME, atribut segmen, dan hak akses. Definisi formal pada subbab ini menjadi sumber kanonik bagi representasi \textit{caveat} Macaroon dan evaluasi pada \textit{Server} PRE.

Elemen \textit{policy} pada DecMed didefinisikan sebagai:

\begin{equation}
    PE_{\text{DecMed}}
    =
    U \cup UA \cup O \cup OA \cup PC,
\end{equation}

dengan [$U$] himpunan pengguna, [$UA$] himpunan atribut pengguna, [$O$] himpunan objek yang dilindungi, [$OA$] himpunan atribut objek, dan [$PC$] himpunan kelas \textit{policy}. Berdasarkan elemen tersebut, model \textit{policy} pada DecMed dirumuskan sebagai:

\begin{equation}
    PM_{\text{DecMed}}
    =
    \langle
    PE_{\text{DecMed}},
    AR,
    ASSIGN,
    ASSOCIATION,
    \mathcal{C}
    \rangle,
\end{equation}

dengan $AR$ sebagai himpunan hak akses, $ASSIGN$ sebagai relasi penugasan, $ASSOCIATION$ sebagai relasi asosiasi hak akses, dan $\mathcal{C}$ sebagai himpunan kondisi akses DecMed.

Elemen $\mathcal{C}$ merupakan perluasan pada rancangan ini untuk merepresentasikan kondisi dinamis yang ditegakkan melalui Macaroon dan \textit{Server} PRE, seperti masa berlaku, rantai delegasi, revokasi, dan bukti kepemilikan identitas aktor pada rantai delegasi.

\subsubsection{Pengguna dan Atribut Pengguna}

Himpunan pengguna $U$ berisi pengguna DecMed yaitu:

\begin{equation}
    U =
    {
    u_{\text{pasien}},
    u_{\text{petugas-rm}},
    u_{\text{perawat}},
    u_{\text{dokter}},
    u_{\text{laboratorium}},
    u_{\text{apoteker}}
    }.
\end{equation}

Elemen pada himpunan tersebut merepresentasikan alamat dompet aktual, bukan hanya nama peran. Sebagai contoh, $u_{\text{dokter-A}}$ dan $u_{\text{dokter-B}}$ merupakan dua pengguna yang berbeda meskipun keduanya memiliki atribut \textit{subrole} dokter. Atribut pengguna dibagi menjadi atribut kelembagaan dan atribut kapabilitas:

\begin{equation}
    UA = UA_{\text{role}} \cup UA_{\text{capability}}.
\end{equation}

Atribut kelembagaan menyatakan \textit{role} atau \textit{subrole} pengguna yang tercatat pada penyimpanan \textit{on-chain}:

\begin{equation}
    UA_{\text{role}} =
        {
            ua_{\text{patient}},
            ua_{\text{administrative-personnel}},
            ua_{\text{nurse}},
            ua_{\text{doctor}},
            ua_{\text{laboratory}},
            ua_{\text{pharmacist}}
        }.
\end{equation}

Relasi penugasan pengguna terhadap atribut perannya dinyatakan sebagai:

\begin{equation}
    u ; ASSIGN ; ua_{\operatorname{role}(u)}
\end{equation}

Atribut kapabilitas merupakan abstraksi terhadap kapabilitas efektif yang diperoleh dari suatu token Macaroon. Untuk token $\tau$, atribut kapabilitasnya dinyatakan sebagai:

\begin{equation}
    ua_{\operatorname{cap}(\tau)}
    \in UA_{\text{capability}}
\end{equation}

Atribut kapabilitas dibentuk secara konseptual ketika \textit{Server} PRE mengevaluasi token dan membentuk himpunan subjek dari \texttt{root\_subject} serta seluruh \texttt{delegated\_to} pada rantai delegasi. Pengguna hanya dianggap memiliki atribut kapabilitas apabila pengguna tersebut termasuk dalam himpunan subjek token dan dapat membuktikan kepemilikan \textit{wallet}:

\begin{equation}
    u ; ASSIGN ; ua_{\operatorname{cap}(\tau)}
    \iff
    \operatorname{ContainsSubject}(\tau,u)
    \land
    \operatorname{WalletProofValid}(u).
\end{equation}

Pemisahan atribut peran dan atribut kapabilitas diperlukan karena keduanya memiliki fungsi yang berbeda. Atribut kapabilitas menentukan cakupan data yang diberikan melalui token, sedangkan atribut peran menentukan apakah pengguna masih memiliki kewenangan kelembagaan yang sesuai. Dengan demikian, kepemilikan \textit{subrole} dokter tidak secara otomatis memberikan akses terhadap seluruh RME pasien.

\subsubsection{Objek dan Atribut Objek}

Himpunan objek $O$ terdiri atas segmen RME terenkripsi yang disimpan pada IPFS:

\begin{equation}
    O = {o_1,o_2,\ldots,o_n}.
\end{equation}

Setiap objek merepresentasikan satu segmen klinis. Informasi mengenai objek, seperti pasien, episode RME, kategori dataset, kategori fungsi, dan CID data terenkripsi, direpresentasikan melalui metadata pada IOTA.

Atribut objek $OA$ digunakan untuk mengelompokkan segmen berdasarkan beberapa tingkat yaitu:

\begin{enumerate}
    \item rumah sakit;
    \item pasien;
    \item episode RME;
    \item kategori dataset; dan
    \item kategori fungsi klinis.
\end{enumerate}

Contoh relasi penugasan untuk satu segmen diagnosis adalah:

\begin{equation}
    \begin{aligned}
        o_{\text{diagnosis-17}}
         & ; ASSIGN ;
        oa_{\text{diagnosis-RME-17}}   \\
         & ; ASSIGN ;
        oa_{\text{rawat-jalan-RME-17}} \\
         & ; ASSIGN ;
        oa_{\text{RME-17}}             \\
         & ; ASSIGN ;
        oa_{\text{pasien-P}}           \\
         & ; ASSIGN ;
        oa_{\text{rumah-sakit-H}}      \\
         & ; ASSIGN ;
        pc_{\text{DecMed}}.
    \end{aligned}
\end{equation}

Himpunan kategori dataset yang digunakan pada rancangan ini adalah $\texttt{RAWAT\_JALAN}$, $\texttt{RAWAT\_INAP}$, $\texttt{LABORATORIUM}$, dan $\texttt{APOTEK}$. Sementara itu, kategori fungsi mengikuti segmentasi data yang telah dijelaskan pada Subbab \ref{subsec:rancangan-segmentasi-data-rme}.

\subsubsection{Hak Akses dan Asosiasi}

Himpunan hak akses yang digunakan adalah:

\begin{equation}
    AR = {\texttt{Read},\texttt{Update}}.
\end{equation}

Hak \texttt{Read} menyatakan kemampuan membaca segmen RME, sedangkan hak \texttt{Update} menyatakan kemampuan menulis atau memperbarui segmen. Pada implementasi token, hak tersebut direpresentasikan melalui \textit{caveat} \texttt{purpose}. Cakupan objeknya direpresentasikan melalui \texttt{read\_dataset\_in}, \texttt{write\_dataset\_in}, \texttt{read\_function\_in}, dan \texttt{write\_function\_in}.

Pada \textit{Policy Machine}, asosiasi dinyatakan sebagai hubungan antara atribut pengguna, himpunan hak akses, dan atribut target:

\begin{equation}
    (ua, ars, oa) \in ASSOCIATION.
\end{equation}

Dalam rancangan DecMed, asosiasi dibentuk secara konseptual dari kapabilitas efektif token:

\begin{equation}
    \left(
    ua_{\operatorname{cap}(\tau)},
    AR_{\operatorname{eff}}(\tau),
    OA_{\operatorname{scope}}(\tau)
    \right)
    \in ASSOCIATION.
\end{equation}

$AR_{\operatorname{eff}}(\tau)$ merupakan hak akses efektif token, sedangkan $OA_{\operatorname{scope}}(\tau)$ merupakan himpunan atribut objek yang termasuk dalam cakupan token. Cakupan objek token didefinisikan sebagai:

\begin{equation}
    \begin{split}
        \operatorname{Scope}(\tau)
        =
        \{ o_s \in O \mid {} &
        p(s)=\operatorname{Patient}(\tau)                                 \\
                             & \land h(s)=\operatorname{Hospital}(\tau)   \\
                             & \land d(s)\in D_{\operatorname{eff}}(\tau) \\
                             & \land f(s)\in F_{\operatorname{eff}}(\tau) \\
                             & \land
        (
        \operatorname{Purpose}(\tau)=\texttt{Update}
        \Rightarrow
        r(s)=\operatorname{RME}(\tau)
        )
        \}.
    \end{split}
\end{equation}

$D_{\operatorname{eff}}(\tau)$ merupakan irisan seluruh \textit{caveats} \textit{dataset category} pada token, sedangkan $F_{\operatorname{eff}}(\tau)$ merupakan irisan seluruh \textit{caveats} \textit{function category}:

\begin{equation}
    D_{\operatorname{eff}}(\tau)
    =
    \bigcap_{i=1}^{n} D_i,
\end{equation}

\begin{equation}
    F_{\operatorname{eff}}(\tau)
    =
    \bigcap_{i=1}^{m} F_i.
\end{equation}

Penggunaan operasi irisan memastikan bahwa penambahan \textit{caveat} pada token turunan hanya dapat mempertahankan atau mempersempit cakupan akses. Penambahan \textit{caveat} baru tidak dapat mengembalikan \textit{dataset category} atau \textit{function category} yang telah dikecualikan oleh token induk.

Untuk operasi \texttt{Read}, token dapat digunakan terhadap seluruh episode RME pasien yang termasuk dalam kategori dataset dan fungsi yang diizinkan. Untuk operasi \texttt{Update}, cakupan juga harus sesuai dengan \texttt{related\_rme\_id} sehingga perubahan hanya dapat dilakukan terhadap episode pelayanan yang sedang menjadi konteks token.

\subsubsection{Kondisi \textit{Policy} Akses}

Selain asosiasi antara kapabilitas dan objek, keputusan akses juga dipengaruhi oleh himpunan kondisi $\mathcal{C}$. Kondisi yang digunakan pada DecMed ditunjukkan pada Tabel \ref{tab:pemetaan-kondisi-policy-decmed}.

\begin{table}[!ht]
    \centering
    \caption{Kondisi \textit{Policy} Akses DecMed}
    \label{tab:pemetaan-kondisi-policy-decmed}
    \small
    \begin{tabular}{|p{0.24\textwidth}|p{0.32\textwidth}|p{0.31\textwidth}|}
        \hline
        \textbf{Kondisi}
         & \textbf{Representasi}
         & \textbf{Fungsi}                                                            \\ \hline

        Identitas pasien
         & \texttt{patient\_address}
         & Membatasi token pada data milik pasien tertentu.                           \\ \hline

        Rumah sakit
         & \texttt{hospital\_cid}
         & Membatasi penggunaan token pada konteks rumah sakit tertentu.              \\ \hline

        Episode RME
         & \texttt{related\_rme\_id}
         & Membatasi operasi \texttt{Update} pada episode RME tertentu.               \\ \hline

        Subjek rantai delegasi
         & \texttt{root\_subject} dan seluruh \texttt{delegated\_to}
         & Menentukan himpunan alamat dompet yang berwenang menggunakan token.        \\ \hline

        Identitas aktor pemohon
         & \textit{Wallet proof}
         & Membuktikan bahwa permintaan ditandatangani oleh salah satu subjek rantai. \\ \hline

        Masa berlaku
         & \texttt{expires\_before}
         & Membatasi waktu penggunaan token.                                          \\ \hline

        Kedalaman delegasi
         & \texttt{max\_delegation\_depth} dan rantai delegasi
         & Membatasi jumlah penurunan token.                                          \\ \hline

        Hubungan token
         & \texttt{parent\_token\_hash}
         & Menghubungkan token turunan dengan token induk untuk revokasi.             \\ \hline

        Otoritas kelembagaan
         & \textit{Role} dan \textit{subrole} pada IOTA
         & Membatasi operasi tulis sesuai kewenangan tenaga kesehatan.                \\ \hline

        Status revokasi
         & Catatan revokasi
         & Menolak token awal atau token turunan yang telah dicabut.                  \\ \hline
    \end{tabular}

\end{table}

Masa berlaku efektif token merupakan nilai paling awal dari seluruh \textit{caveat} \texttt{expires\_before}:

\begin{equation}
    Expiry_{\operatorname{eff}}(\tau)
    =
    \min
    {
        Expiry_1,Expiry_2,\ldots,Expiry_n
    }.
\end{equation}

Himpunan subjek yang berwenang menggunakan token ditentukan berdasarkan rantai delegasi:

\begin{equation}
    \operatorname{Subjects}(\tau)
    =
    \left\{
    \operatorname{root\_subject}(\tau)
    \right\}
    \cup
    \left\{
    \operatorname{delegated\_to}_i(\tau)
    \mid
    1 \leq i \leq n
    \right\},
\end{equation}

dan permintaan dari aktor $u$ hanya diterima apabila:

\begin{equation}
    \operatorname{ContainsSubject}(\tau,u)
    \iff
    u \in \operatorname{Subjects}(\tau).
\end{equation}

Kedalaman delegasi dihitung berdasarkan jumlah perpindahan aktor pada rantai delegasi dan harus memenuhi:

\begin{equation}
    Depth(\tau)
    \leq
    MaxDepth(\tau).
\end{equation}

\subsubsection{Rancangan \textit{Policy} Delegasi}

Delegasi dimodelkan sebagai pembentukan kapabilitas turunan $\tau_c$ dari kapabilitas induk $\tau_p$. Token turunan hanya valid apabila tidak memperluas hak akses token induk.

Hubungan cakupan token harus memenuhi:

\begin{equation}
    \operatorname{Scope}(\tau_c)
    \subseteq
    \operatorname{Scope}(\tau_p).
\end{equation}

Hak akses token turunan juga harus memenuhi:

\begin{equation}
    AR_{\operatorname{eff}}(\tau_c)
    \subseteq
    AR_{\operatorname{eff}}(\tau_p).
\end{equation}

Batas waktu token turunan tidak boleh melebihi batas waktu token induk:

\begin{equation}
    Expiry_{\operatorname{eff}}(\tau_c)
    \leq
    Expiry_{\operatorname{eff}}(\tau_p).
\end{equation}

Selain itu, token turunan harus tetap berada dalam konteks pasien dan rumah sakit yang sama:

\begin{equation}
    Patient(\tau_c)=Patient(\tau_p),
\end{equation}

\begin{equation}
    Hospital(\tau_c)=Hospital(\tau_p).
\end{equation}

Untuk token dengan operasi \texttt{Update}, episode RME token turunan juga tidak boleh berbeda dari token induk:

\begin{equation}
    RME(\tau_c)=RME(\tau_p).
\end{equation}

Hubungan antara token turunan dan token induk disimpan melalui \texttt{parent\_token\_hash}. Sementara itu, perubahan aktor direpresentasikan dengan menambahkan \texttt{delegated\_by} dan \texttt{delegated\_to}. Karena seluruh \textit{caveat} token induk tetap dipertahankan, proses delegasi hanya dapat menambahkan \textit{caveats} baru dan tidak dapat menghapus \textit{caveats} sebelumnya.

\subsubsection{Fungsi Keputusan Akses}

Suatu permintaan akses terhadap segmen RME direpresentasikan sebagai:

\begin{equation}
    q =
    \langle
    u,h,p,r,d,f,a,t,\tau,\pi
    \rangle,
    \label{eq:permintaan-akses}
\end{equation}

dengan [$u$] alamat dompet pengguna yang melakukan permintaan, [$h$] identitas rumah sakit, [$p$] alamat pasien, [$r$] identitas episode RME, [$d$] kategori dataset, [$f$] kategori fungsi, [$a$] operasi yang diminta, [$t$] waktu permintaan, [$\tau$] token Macaroon, dan [$\pi$] bukti kepemilikan \textit{wallet}. Permintaan diberikan apabila seluruh kondisi \textit{policy} terpenuhi yaitu:

\begin{equation}
    \begin{split}
        \operatorname{Permit}(q)
        \iff {} &
        \operatorname{TokenValid}(\tau)                           \\
                & \land \operatorname{PurposeMatch}(a,\tau)       \\
                & \land \operatorname{ContainsSubject}(\tau,u)    \\
                & \land \operatorname{WalletProofValid}(\pi,q,u)  \\
                & \land \operatorname{PatientMatch}(p,\tau)       \\
                & \land \operatorname{HospitalMatch}(h,\tau)      \\
                & \land \operatorname{DatasetAllowed}(d,\tau)     \\
                & \land \operatorname{FunctionAllowed}(f,\tau)    \\
                & \land \operatorname{TimeValid}(t,\tau)          \\
                & \land \operatorname{DelegationChainValid}(\tau) \\
                & \land \neg\operatorname{Revoked}(\tau)
    \end{split}
    \label{eq:permit-akses}
\end{equation}

Apabila salah satu kondisi tidak terpenuhi, keputusan akses adalah:

\begin{equation}
    \operatorname{Decision}(q)=\operatorname{Deny}.
    \label{eq:deny-akses}
\end{equation}

Prinsip \textit{deny by default} diterapkan melalui fungsi tersebut. Kepemilikan token saja tidak cukup untuk memperoleh akses. Token harus valid, permintaan harus berasal dari salah satu aktor pada rantai delegasi, objek harus berada dalam cakupan token, token tidak boleh kedaluwarsa atau direvokasi, dan operasi tulis harus sesuai dengan \textit{subrole} tenaga kesehatan.

\subsubsection{Pemetaan Model \textit{Policy} ke Komponen DecMed}

Pemetaan elemen model \textit{policy} terhadap komponen implementasi ditunjukkan pada Tabel \ref{tab:pemetaan-policy-machine-decmed}.

\begin{table}[!ht]
    \centering
    \caption{Pemetaan Model \textit{Policy} terhadap Komponen DecMed}
    \label{tab:pemetaan-policy-machine-decmed}
    \small
    \begin{tabular}{|p{0.18\textwidth}|p{0.29\textwidth}|p{0.39\textwidth}|}
        \hline
        \textbf{Elemen}
         & \textbf{Representasi DecMed}
         & \textbf{Sumber atau Mekanisme Implementasi}                              \\ \hline

        $U$
         & Pengguna berdasarkan \textit{wallet address}
         & \textit{wallet address} dan tanda tangan \textit{wallet}.                \\ \hline

        $UA_{\text{role}}$
         & \textit{Role} dan \textit{subrole} pengguna
         & \textit{Registry Personnel} pada IOTA.                                   \\ \hline

        $UA_{\text{capability}}$
         & Kapabilitas efektif pengguna
         & Token Macaroon dan rantai delegasi.                                      \\ \hline

        $O$
         & Segmen RME terenkripsi
         & Objek data pada IPFS.                                                    \\ \hline

        $OA$
         & Rumah sakit, pasien, RME, dataset, dan fungsi
         & Metadata segmen pada IOTA.                                               \\ \hline

        $PC$
         & Kelas \textit{policy} DecMed
         & Aturan kontrol akses yang diberlakukan pada sistem DecMed.               \\ \hline

        $AR$
         & Operasi \texttt{Read} dan \texttt{Update}
         & \texttt{purpose} dan \textit{caveats} baca atau tulis pada Macaroon.     \\ \hline

        $ASSIGN$
         & Penugasan pengguna ke peran dan segmen ke kategori
         & Registri personel serta struktur metadata segmen.                        \\ \hline

        $ASSOCIATION$
         & Hubungan kapabilitas dengan hak dan cakupan objek
         & \textit{Effective capability} hasil evaluasi \textit{caveat}.            \\ \hline

        $\mathcal{C}$
         & Masa berlaku, delegasi, revokasi, dan bukti identitas
         & Macaroon, \textit{wallet proof}, Redis, IOTA, dan fungsi verifikasi PRE. \\ \hline
    \end{tabular}

\end{table}

Berdasarkan rancangan tersebut, model CapBAC yang telah digunakan oleh DecMed tidak digantikan oleh model PBAC. Macaroon tetap berfungsi sebagai token kapabilitas yang dibawa oleh pengguna, sedangkan elemen dan kondisi yang harus diperiksa untuk menafsirkan kapabilitas tersebut ditentukan oleh model PBAC. Dengan demikian, cara otoritas dibawa oleh token dijelaskan melalui CapBAC, sedangkan \textit{policy} yang menentukan validitas penggunaan otoritas tersebut dijelaskan melalui PBAC.

Rincian representasi \textit{policy} dalam bentuk \textit{caveat} dijelaskan pada Subbab \ref{subsec:rancangan-token-macaroons}. Mekanisme evaluasi dan penegakannya pada \textit{Server} PRE dijelaskan pada Subbab \ref{subsec:rancangan-penegakan-otoritas}.
